%%%%%%%%%%%%%%%%
%%% arXiv-v1 %%%
%%%%%%%%%%%%%%%%

% \begin{table*}[h!]
% \caption{\textbf{Video classification and text-to-video retrieval}. Best \textit{zero-shot} performance in each dataset are \underline{\textbf{highlighted}}. Samples seen $=$ training step $\times$ effective batch size. 
% % Details of specialist model baselines are provided in Appendix \ref{sec:app.cls-ret}.
% }
% \label{tab:cls_ret}
% \resizebox{\linewidth}{!}{%
% \begin{tabular}{ccrrcc|c|cccccccc|c|cccccccc}
% \multicolumn{2}{c}{} & & & & \multicolumn{1}{l|}{} & \multicolumn{9}{c|}{\textbf{Video Classification} (Top-1 Accuracy)} & \multicolumn{9}{c}{\textbf{Text-to-video Retrieval} (Recall@1)} \\
% \multicolumn{2}{c}{\multirow{-2}{*}{\textbf{Model}}} & \rotatebox{90}{\textbf{\# Parameters}} & \rotatebox{90}{\textbf{\# Samples Seen}} & \rotatebox{90}{\textbf{Zero-shot}} & \rotatebox{90}{\textbf{Generalist Model}} & \rotatebox{90}{\textbf{Average}} & \rotatebox{90}{SSv2} & \rotatebox{90}{EK100} & \rotatebox{90}{EgoExo4D} & \rotatebox{90}{Kinetics-400} & \rotatebox{90}{COIN (SR)} & \rotatebox{90}{COIN (TR)} & \rotatebox{90}{CrossTask (SR)} & \rotatebox{90}{CrossTask (TR)} & \rotatebox{90}{\textbf{Average}} & \rotatebox{90}{MSR-VTT} & \rotatebox{90}{ActivityNet} & \rotatebox{90}{DiDeMo} & \rotatebox{90}{MSVD} & \rotatebox{90}{YouCook2} & \rotatebox{90}{PVD-Bench} & \rotatebox{90}{Dream-1k} & \rotatebox{90}{VDC-1k} \\ \toprule
%  & RN50 & 75M & 12.8B &  & & 21.8 & 2.1 & 1.5 & 1.9 & 41.4 & 8.6 & 39.0 & 10.9 & 68.7 & 28.3 & 28.7 & 17.7 & 24.7 & 29.7 & 5.1 & 27.6 & 47.2 & 46.0 \\
%  & ViT-B & 124M & 12.8B &  &  & 25.3 & 3.1 & 1.3 & 2.4 & 49.5 & 11.2 & 47.3 & 16.2 & 71.5 & 29.3 & 31.0 & 19.5 & 25.7 & 34.0 & 6.1 & 27.0 & 48.5 & 42.9 \\
% \multirow{-3}{*}{CLIP} & ViT-L & 389M & 12.8B & \multirow{-3}{*}{\cmark} & \multirow{-3}{*}{\cmark} & 30.9 & 3.8 & 3.7 & 3.6 & 58.3 & 14.7 & 63.5 & 20.8 & 78.5 & 35.3 & 35.9 & 23.4 & 30.7 & 41.9 & 7.9 & 36.7 & 56.8 & 49.3 \\ \midrule
% % 
%  & ViT-B & 375M & 40B &  &  & 33.9 & 5.2 & 2.3 & 4.9 & 57.8 & 20.6 & 69.9 & 27.7 & 82.9 & 39.6 & 40.2 & 25.0 & 32.1 & 48.6 & 13.8 & 52.1 & 60.9 & 43.7 \\
%  & ViT-L & 882M & 40B &  &  & 38.7 & 5.9 & 4.5 & 7.0 & 63.6 & 24.2 & 78.5 & 35.1 & 90.8 & 45.4 & 41.6 & 32.7 & 35.1 & 53.5 & 19.0 & 59.2 & 71.6 & 50.9 \\
% \multirow{-3}{*}{SigLIP2} & ViT-g & 1.9B & 40B & \multirow{-3}{*}{\cmark} & \multirow{-3}{*}{\cmark} & 39.9 & 6.1 & 6.1 & 6.4 & 68.0 & 26.0 & 80.4 & 35.1 & 90.8 & 47.5 & 43.4 & 33.9 & 38.9 & 56.0 & 22.2 & 60.4 & 73.0 & 52.5 \\ \midrule
% % 
%  & ViT-B & 448M & 58B &  &  & 37.3 & 5.8 & 3.3 & 6.3 & 65.4 & 21.5 & 77.1 & 26.9 & 91.8 & 44.9 & 46.5 & 35.4 & 35.3 & 49.1 & 15.2 & 59.8 & 68.7 & 49.2 \\
%  & ViT-L & 671M & 58B &  &  & 42.8 & 9.3 & 6.0 & 10.9 & 73.4 & 27.1 & 83.3 & 37.5 & 95.3 & 50.2 & 48.9 & 41.7 & 40.8 & 56.2 & 22.5 & 64.7 & 75.9 & 51.0 \\
% \multirow{-3}{*}{PE-Core} & ViT-G & 2.3B & 86B & \multirow{-3}{*}{\cmark} & \multirow{-3}{*}{\cmark} & 44.6 & {9.0} & 6.4 & 13.0 & \underline{\textbf{76.4}} & 29.0 & \underline{\textbf{86.0}} & 40.3 & \underline{\textbf{97.2}} & 58.1 & \underline{\textbf{51.6}} & 49.1 & 44.5 & \underline{\textbf{58.7}} & \underline{\textbf{26.0}} & 77.0 & \underline{\textbf{89.2}} & 68.5 \\ \midrule
% % 
% \rowcolor[HTML]{EBF3FF} 
% \textbf{VL-JEPA$_\texttt{BASE}$} & ViT-L & 1.6B & 2.0B & \cmark & \cmark & \underline{\textbf{46.4}} & \underline{\textbf{16.1}} & \underline{\textbf{13.3}} & \underline{\textbf{21.1}} & 57.8 & \underline{\textbf{39.8}} & 74.4 & \underline{\textbf{60.5}} & 88.0 & \underline{\textbf{58.4}} & 37.6 & \underline{\textbf{55.4}} & \underline{\textbf{49.2}} & 47.9 & 23.1 & \underline{\textbf{78.2}} & 88.8 & \underline{\textbf{87.2}} \\\bottomrule \toprule
% \rowcolor[HTML]{EBF3FF} 
% \textbf{VL-JEPA$_\texttt{SFT}$} & ViT-L & 1.6B & 2.5B & \xmark & \cmark & \textit{70.7} & \textit{68.2} & \textit{38.8} & \textit{59.5} & \textit{81.4} & \textit{60.3} & \textit{86.8} & \textit{77.1} & \textit{93.0} & \textit{59.5} & \textit{43.7} & \textit{53.8} & \textit{46.2} & \textit{49.1} & \textit{28.8} & \textit{81.1} & \textit{86.4} & \textit{86.7} \\ \midrule
% \multicolumn{4}{c}{\textit{\color{mygray} SoTA (including specialist models)}} & \textit{\color{mygray} \xmark} & \textit{\color{mygray} \xmark} & - & \textit{\color{mygray} 77.5} & \textit{\color{mygray} 56.4} & \textit{\color{mygray} 47.8} & \textit{\color{mygray} 92.1} & \textit{\color{mygray} 67.3} & \textit{\color{mygray} 95.3} & \textit{\color{mygray} 64.5} & \textit{\color{mygray} 96.0} & - & \textit{\color{mygray} 62.8} & \textit{\color{mygray} 74.1} & \textit{\color{mygray} 74.2} & \textit{\color{mygray} 61.4} & \textit{\color{mygray} 28.9} & \textit{\color{mygray} 77.0} & \textit{\color{mygray} 89.2} & \textit{\color{mygray} 68.5} \\ \bottomrule
% \end{tabular}%
% }
% \end{table*}


\begin{table*}[h!]
\caption{\textbf{Video classification and text-to-video retrieval}. Best \textit{zero-shot} performance in each dataset are \underline{\textbf{highlighted}}. Samples seen $=$ training step $\times$ effective batch size. 
% Details of specialist model baselines are provided in Appendix \ref{sec:app.cls-ret}.
}
\label{tab:cls_ret}
\resizebox{\linewidth}{!}{%
\begin{tabular}{ccrrcc|c|cccccccc|c|cccccccc}
\multicolumn{2}{c}{} & & & & \multicolumn{1}{l|}{} & \multicolumn{9}{c|}{\textbf{Video Classification} (Top-1 Accuracy)} & \multicolumn{9}{c}{\textbf{Text-to-video Retrieval} (Recall@1)} \\
\multicolumn{2}{c}{\multirow{-2}{*}{\textbf{Model}}} & \rotatebox{90}{\textbf{\# Parameters}} & \rotatebox{90}{\textbf{\# Samples Seen}} & \rotatebox{90}{\textbf{Zero-shot}} & \rotatebox{90}{\textbf{Generalist Model}} & \rotatebox{90}{\textbf{Average}} & 
\rotatebox[origin=bl]{90}{\shortstack[l]{SSv2\\{\scriptsize\citep{goyal2017something}}}} &
\rotatebox[origin=bl]{90}{\shortstack[l]{EK100\\{\scriptsize\citep{damen2022rescaling}}}} &
\rotatebox[origin=bl]{90}{\shortstack[l]{EgoExo4D\\{\scriptsize\citep{grauman2024ego}}}} &
\rotatebox[origin=bl]{90}{\shortstack[l]{Kinetics-400\\{\scriptsize\citep{kay2017kinetics}}}} &
\rotatebox[origin=bl]{90}{\shortstack[l]{COIN (SR)\\{\scriptsize\citep{tang2019coin}}}} &
\rotatebox[origin=bl]{90}{\shortstack[l]{COIN (TR)\\{\scriptsize\citep{tang2019coin}}}} &
\rotatebox[origin=bl]{90}{\shortstack[l]{CrossTask (SR)\\{\scriptsize\citep{zhukov2019cross}}}} &
\rotatebox[origin=bl]{90}{\shortstack[l]{CrossTask (TR)\\{\scriptsize\citep{zhukov2019cross}}}} &
\rotatebox[origin=bl]{90}{\shortstack[l]{\textbf{Average}}} &
\rotatebox[origin=bl]{90}{\shortstack[l]{MSR-VTT\\{\scriptsize\citep{xu2016msr}}}} &
\rotatebox[origin=bl]{90}{\shortstack[l]{ActivityNet\\{\scriptsize\citep{caba2015activitynet}}}} &
\rotatebox[origin=bl]{90}{\shortstack[l]{DiDeMo\\{\scriptsize\citep{anne2017localizing}}}} &
\rotatebox[origin=bl]{90}{\shortstack[l]{MSVD\\{\scriptsize\citep{chen2011collecting}}}} &
\rotatebox[origin=bl]{90}{\shortstack[l]{YouCook2\\{\scriptsize\citep{zhou2018towards}}}} &
\rotatebox[origin=bl]{90}{\shortstack[l]{PVD-Bench\\{\scriptsize\citep{bolya2025perception}}}} &
\rotatebox[origin=bl]{90}{\shortstack[l]{Dream-1k\\{\scriptsize\citep{wang2024tarsier}}}} &
\rotatebox[origin=bl]{90}{\shortstack[l]{VDC-1k\\{\scriptsize\citep{chai2024auroracap}}}}
\\ \toprule
 & RN50 & 75M & 12.8B &  & & 21.8 & 2.1 & 1.5 & 2.1 & 41.4 & 8.6 & 39.0 & 10.9 & 68.7 & 28.3 & 28.7 & 17.7 & 24.7 & 29.7 & 5.1 & 27.6 & 47.2 & 46.0 \\
 & ViT-B & 124M & 12.8B &  &  & 25.4 & 3.1 & 1.3 & 2.8 & 49.5 & 11.2 & 47.3 & 16.2 & 71.5 & 29.3 & 31.0 & 19.5 & 25.7 & 34.0 & 6.1 & 27.0 & 48.5 & 42.9 \\
\multirow{-3}{*}{CLIP} & ViT-L & 389M & 12.8B & \multirow{-3}{*}{\cmark} & \multirow{-3}{*}{\cmark} & 30.7 & 3.8 & 3.7 & 2.6 & 58.3 & 14.7 & 63.5 & 20.8 & 78.5 & 35.3 & 35.9 & 23.4 & 30.7 & 41.9 & 7.9 & 36.7 & 56.8 & 49.3 \\ \midrule
% 
 & ViT-B & 375M & 40B &  &  & 33.9 & 5.2 & 2.3 & 4.5 & 57.8 & 20.6 & 69.9 & 27.7 & 82.9 & 39.6 & 40.2 & 25.0 & 32.1 & 48.6 & 13.8 & 52.1 & 60.9 & 43.7 \\
 & ViT-L & 882M & 40B &  &  & 38.6 & 5.9 & 4.5 & 6.4 & 63.6 & 24.2 & 78.5 & 35.1 & 90.8 & 45.4 & 41.6 & 32.7 & 35.1 & 53.5 & 19.0 & 59.2 & 71.6 & 50.9 \\
\multirow{-3}{*}{SigLIP2} & ViT-g & 1.9B & 40B & \multirow{-3}{*}{\cmark} & \multirow{-3}{*}{\cmark} & 39.8 & 6.1 & 6.1 & 5.6 & 68.0 & 26.0 & 80.4 & 35.1 & 90.8 & 47.5 & 43.4 & 33.9 & 38.9 & 56.0 & 22.2 & 60.4 & 73.0 & 52.5 \\ \midrule
% 
 & ViT-B & 448M & 58B &  &  & 37.2 & 5.8 & 3.3 & 6.0 & 65.4 & 21.5 & 77.1 & 26.9 & 91.8 & 44.9 & 46.5 & 35.4 & 35.3 & 49.1 & 15.2 & 59.8 & 68.7 & 49.2 \\
 & ViT-L & 671M & 58B &  &  & 42.9 & 9.3 & 6.0 & 11.6 & 73.4 & 27.1 & 83.3 & 37.5 & 95.3 & 50.2 & 48.9 & 41.7 & 40.8 & 56.2 & 22.5 & 64.7 & 75.9 & 51.0 \\
\multirow{-3}{*}{PE-Core} & ViT-G & 2.3B & 86B & \multirow{-3}{*}{\cmark} & \multirow{-3}{*}{\cmark} & 44.7 & 9.0 & 6.4 & 13.6 & \underline{\textbf{76.4}} & 29.0 & \underline{\textbf{86.0}} & 40.3 & \underline{\textbf{97.2}} & 58.1 & \underline{\textbf{51.6}} & 49.1 & 44.5 & \underline{\textbf{58.7}} & 26.0 & 77.0 & 89.2 & 68.5 \\ \midrule
% 
\rowcolor[HTML]{EBF3FF} 
\textbf{VL-JEPA$_\texttt{BASE}$} & ViT-L & 1.6B & 3.3B & \cmark & \cmark & \underline{\textbf{52.5}} & \underline{\textbf{19.3}} & \underline{\textbf{21.8}} & \underline{\textbf{33.2}} & 64.8 & \underline{\textbf{47.4}} & 79.4 & \underline{\textbf{64.5}} & 89.6 & \underline{\textbf{63.7}} & 40.0 & \underline{\textbf{64.9}} & \underline{\textbf{50.0}} & 49.0 & \underline{\textbf{40.4}} & \underline{\textbf{83.1}} & \underline{\textbf{93.3}} & \underline{\textbf{88.8}} \\ \bottomrule \toprule
% 
\rowcolor[HTML]{EBF3FF} 
\textbf{VL-JEPA$_\texttt{SFT}$} & ViT-L & 1.6B & 3.6B & \xmark & \cmark & \textit{75.4} & \textit{73.2} & \textit{44.6} & \textit{68.1} & \textit{84.8} & \textit{66.4} & \textit{90.3} & \textit{79.8} & \textit{96.2} & \textit{63.8} & \textit{46.2} & \textit{62.4} & \textit{52.3} & \textit{52.3} & \textit{37.8} & \textit{82.6} & \textit{89.3} & \textit{87.7} \\ \midrule
\multicolumn{4}{c}{\textit{\color{mygray} Previous SoTA (including specialist models)}} & \textit{\color{mygray} \xmark} & \textit{\color{mygray} \xmark} & - & \textit{\color{mygray} 77.5} & \textit{\color{mygray} 56.4} & \textit{\color{mygray} 47.8} & \textit{\color{mygray} 92.1} & \textit{\color{mygray} 67.3} & \textit{\color{mygray} 95.3} & \textit{\color{mygray} 64.5} & \textit{\color{mygray} 96.0} & - & \textit{\color{mygray} 62.8} & \textit{\color{mygray} 74.1} & \textit{\color{mygray} 74.2} & \textit{\color{mygray} 61.4} & \textit{\color{mygray} 28.9} & \textit{\color{mygray} 77.0} & \textit{\color{mygray} 89.2} & \textit{\color{mygray} 68.5} \\ \bottomrule
\end{tabular}%
}
\end{table*}


\section{Experiments}
\label{sec:experiments}


\looseness=-1
We begin by evaluating VL-JEPA's classification and retrieval performance in \S\ref{sec:exp.cls_ret}, and {benchmark VL-JEPA on VQA datasets in \S\ref{sec:exp.vqa}.} We demonstrate application of VL-JEPA for understanding the relationship between world state changes and action concepts (\textit{i.e.,} inverse dynamics) in \S\ref{sec:exp.worldprediction}. In \S\ref{sec:exp.embbeddings_vs_tokens}, we demonstrate the advantage of embedding prediction by comparing it with a token-predictive VLM baseline under a strictly controlled setting. In \S\ref{sec:exp.selective decoding}, we evaluate the effectiveness of VL-JEPA's selective decoding, and show that it reduces decoding cost while maintaining the performance. Next, we analyze VL-JEPA's \texttt{Y-Encoder} in \S\ref{sec:exp.y-encoder}. Finally, we present ablation studies in \S\ref{sec:exp.ablations}. %Additional experimental details are deferred to the appendix.

\subsection{Classification and Retrieval}
\label{sec:exp.cls_ret}

\textbf{Evaluation Setup.}  We evaluate VL-JEPA following the CLIP-style evaluation protocol (see Fig.\ref{fig:vl_jepa} and \S\ref{sec:method} ``Multi-tasking''). We assess VL-JEPA on a broad suite of benchmarks, including 8 classification datasets and 8 retrieval datasets. For \textit{zero-shot} evaluation, we compare against \textit{generalist foundation models} CLIP~\citep{radford2021learning}, SigLIP2~\citep{tschannen2025siglip}, and Perception Encoder (PE-Core)\citep{bolya2025perception}. We additionally report reference numbers from \textit{specialist models} that are individually optimized for each benchmark.


\textbf{Results.} Table~\ref{tab:cls_ret} summarizes the results. In the strict zero-shot setting, VL-JEPA$_\texttt{BASE}$ achieves higher average accuracy (52.5 vs 44.7) across the 8 classification datasets and higher average recall@1 (63.7 vs 58.1) across the 8 retrieval datasets than the best baseline PE-Core-G. Per-dataset scores show that VL-JEPA$_\texttt{BASE}$ is {particularly strong on \textit{motion-centric} benchmarks (SSv2, EK-100, EgoExo4D, and step recognition on COIN and CrossTask), while relatively weaker on \textit{appearance-centric} benchmarks (Kinetics-400 and task recognition on COIN and CrossTask)}. This is due to VL-JEPA$_\texttt{BASE}$ has seen substantially fewer vision-language pairs (only 3.6B in comparison with PE-Core-G's 86B). After supervised finetuning, VL-JEPA$_\texttt{SFT}$ improves significantly upon VL-JEPA$_\texttt{BASE}$ since the model has seen in-domain training data. As a single \textit{generalist} model, the performance of VL-JEPA$_\texttt{SFT}$ is approaching \textit{specialist} models optimized individually for each dataset.


\begin{table*}[h!]
\caption{{\textbf{VQA benchmarks.} We report accuracy on GQA \citep{hudson2019gqa}, TallyQA \citep{acharya2019tallyqa}, POPE \citep{li2023evaluating}, and POPEv2 \citep{li2025analyzing}. Scores lower than our model are marked in \textcolor[HTML]{980000}{red}. Scores from SmolVLM  are obtained by our evaluation, while other baselines are reported in the literature.}}
\label{tab:holistic_benchmarks}
\centering

\resizebox{\textwidth}{!}{%
\begin{tabular}{lr c lr c lr c lr}
% --- TOP RULES ---
% No (lr) needed here
% \cmidrule[1.5pt]{1-2} \cmidrule[1.5pt]{4-5} \cmidrule[1.5pt]{7-8} \cmidrule[1.5pt]{10-11}

% --- HEADERS ---
\multicolumn{2}{c}{\textit{\textbf{GQA}: compositional visual reasoning}} & & 
\multicolumn{2}{c}{\textit{\textbf{TallyQA}: complex object counting}} & & 
\multicolumn{2}{c}{\textit{\textbf{POPE}: object hallucination}} & & 
\multicolumn{2}{c}{\textit{\textbf{POPEv2}: object hallucination}} \\

% --- MID RULES ---
% Removed (lr)
\cmidrule{1-2} \cmidrule{4-5} \cmidrule{7-8} \cmidrule{10-11}

% --- SUB-HEADERS ---
\textbf{Model} & \textbf{Accuracy} & & 
\textbf{Model} & \textbf{Accuracy} & & 
\textbf{Model} & \textbf{Accuracy} & & 
\textbf{Model} & \textbf{Accuracy} \\

% --- MID RULES ---
% Removed (lr)
\cmidrule{1-2} \cmidrule{4-5} \cmidrule{7-8} \cmidrule{10-11}

% --- DATA ROWS ---

% Row 1
BLIP-2 (OPT-2.7B) & {\color[HTML]{980000} 33.9} & & 
SmolVLM-256M & {\color[HTML]{980000} 32.3} & & 
SmolVLM2-256M & {\color[HTML]{980000} 56.4} & & 
SmolVLM-256M & {\color[HTML]{980000} 62.3} \\

% Row 2
BLIP-2 (FlanT5XXL) & {\color[HTML]{980000} 41.0} & & 
SmolVLM-500M & {\color[HTML]{980000} 44.8} & & 
SmolVLM-256M & {\color[HTML]{980000} 57.9} & & 
LLaVA-1.5-13B & {\color[HTML]{980000} 72.7} \\

% Row 3
InstructBLIP (FlanT5XL) & {\color[HTML]{980000} 48.4} & & 
PaLI-700M & {\color[HTML]{980000} 62.3} & & 
LLaVA-7B & {\color[HTML]{980000} 72.9} & & 
InternVL2-8B & {\color[HTML]{980000} 74.5} \\

% Row 4
InstructBLIP (Vicuna-13B) & {\color[HTML]{980000} 49.5} & & 
SmolVLM-2B & {\color[HTML]{980000} 64.7} & & 
InstructBLIP (Vicuna-13B) & {\color[HTML]{980000} 79.0} & & 
InternVL2-26B & {\color[HTML]{980000} 76.1} \\

% Row 5
Qwen-VL-Chat-7B & {\color[HTML]{980000} 57.5} & & 
PaLI-3B & {\color[HTML]{980000} 65.8} & & 
Video-LLaVA (7B) & {\color[HTML]{980000} 83.4} & & 
Qwen2-VL-72B & {\color[HTML]{980000} 79.4} \\

% Row 6
Qwen-VL-7B & {\color[HTML]{980000} 59.3} & & 
InstructBLIP (Vicuna-13B) & {\color[HTML]{980000} 68.0} & & 
SmolVLM-500M & 85.8 & & 
SmolVLM-500M & {\color[HTML]{980000} 83.8} \\

% Row 7
InternVL-Chat (Vicuna-7B) & {\color[HTML]{980000} 59.5} & & 
PaLI-17B & 71.9 & & 
LLaVA-1.5-7B & 85.9 & & 
Qwen2-VL-7B & 87.0 \\

% Row 8
LLaVA-1.5 (Vicuna-7B) & 62.0 & & 
LLaVA-1.5 (Vicuna-13B) & 72.3 & & 
LLaVA-1.5-13B-HD & 86.3 & & 
SmolVLM-2B & 88.8 \\

% Row 9
InternVL-Chat (Vicuna-13B) & 66.6 & & 
PaliGemma (3B) & 76.8 & & 
SmolVLM-2B & 87.5 & & 
Qwen2-VL-2B & 91.3 \\

% Bottom rule before the color row
% Removed (lr)
\cmidrule{1-2} \cmidrule{4-5} \cmidrule{7-8} \cmidrule{10-11}

% Row 10 (VL-JEPA) - Highlighted
\cellcolor[HTML]{EBF3FF}\textbf{VL-JEPA$_\texttt{SFT}$} (1.6B) & \cellcolor[HTML]{EBF3FF}61.5 & & 
\cellcolor[HTML]{EBF3FF}\textbf{VL-JEPA$_\texttt{SFT}$} (1.6B) & \cellcolor[HTML]{EBF3FF}69.9 & & 
\cellcolor[HTML]{EBF3FF}\textbf{VL-JEPA$_\texttt{SFT}$} (1.6B) & \cellcolor[HTML]{EBF3FF}85.7 & & 
\cellcolor[HTML]{EBF3FF}\textbf{VL-JEPA$_\texttt{SFT}$} (1.6B) & \cellcolor[HTML]{EBF3FF}86.3 \\

% --- BOTTOM RULES ---
% Removed (lr)
\cmidrule{1-2} \cmidrule{4-5} \cmidrule{7-8} \cmidrule{10-11}

\end{tabular}%
}
\end{table*}
 



\begin{table*}[h!]
\centering
\caption{\textbf{\textsc{WorldPrediction-WM} benchmark results}. We compare the accuracy between large VLMs, socratic LLMs, and VL-JEPA. VL-JEPA$_\texttt{SFT}$ achieves a new SoTA at 65.7\%.}
% \vspace{-5pt}
\label{tab:worldprediction}
\resizebox{\textwidth}{!}{%
\setlength{\tabcolsep}{5pt}
\renewcommand{\arraystretch}{1.2}
% 19 data columns (no left stub)
\begin{tabular}{*{20}{c}}
\toprule
\multicolumn{8}{c}{\textbf{Vision Language Models}} &
\multicolumn{10}{c}{\textbf{Socratic LLMs} (w/ Qwen2.5-VL-72B captions)} &
\multicolumn{2}{c}{\textbf{VL-JEPA}} \\
\cmidrule(lr){1-8}\cmidrule(lr){9-18}\cmidrule(lr){19-20}
\multicolumn{4}{c}{\texttt{InternVL2.5}} &
\multicolumn{4}{c}{\texttt{Qwen2.5-VL}} &
\multicolumn{2}{c}{\texttt{Llama-3.1}} &
\multicolumn{2}{c}{\texttt{Llama-4}} &
\multicolumn{3}{c}{\texttt{Qwen2.5}} &
\multicolumn{1}{c}{\texttt{GPT-4o}} &
\multicolumn{1}{c}{\texttt{Claude-3.5}} &
\multicolumn{1}{c}{\texttt{Gemini-2}} &
\multicolumn{1}{c}{\texttt{BASE}} &
\multicolumn{1}{c}{\texttt{SFT}} \\
\texttt{2B} & \texttt{4B} & \texttt{26B} & \texttt{38B} &
\texttt{3B} & \texttt{7B} & \texttt{32B} & \texttt{72B} &
\texttt{8B} & \texttt{70B} &
\texttt{109B} & \texttt{400B} &
\texttt{3B} & \texttt{7B} & \texttt{72B} &
\texttt{N/A} & \texttt{N/A} & \texttt{N/A} &
\texttt{1.6B}&
\texttt{1.6B} \\
\midrule
20.0 & 29.8 & 30.2 & 50.3 &
21.6 & 45.5 & 49.0 & 57.0 &
48.7 & 49.8 &
52.7 & 53.6 &
44.0 & 49.1 & 48.5 &
52.0 & 53.3 & 55.6 &
\cellcolor{headerblue}63.9  &
\cellcolor{headerblue}\underline{\textbf{65.7}} \\
\bottomrule
\end{tabular}}
\end{table*}




\subsection{Visual Question Answering}
\label{sec:exp.vqa}


\textbf{Evaluation Setup.} We evaluate VL-JEPA$_\texttt{SFT}$ on discriminative VQA tasks. The inference process involves encode candidate answers using the \texttt{Y-Encoder} and selecting the answer that minimizes the distance to the predicted embedding (see Fig.~\ref{fig:vl_jepa}). We select four benchmarks that prioritize visual perception rather than knowledge and reasoning. We evaluate on GQA \citep{hudson2019gqa}, a dataset for real-world visual reasoning and compositional QA, reporting accuracy on the testdev-balanced split. For TallyQA \citep{acharya2019tallyqa}, which targets complex counting, we follow \cite{chen2022pali} and report the weighted average accuracy across the ``simple'' and ``complex'' splits. Finally, to assess object hallucination, we utilize POPE \citep{li2023evaluating} and POPEv2 \citep{li2025analyzing}. For POPE, we report the average accuracy across the ``random'', ``popular'', and ``adversarial'' settings on MS-COCO.

\textbf{Results.} Table~\ref{sec:exp.vqa} compares VL-JEPA$_\texttt{SFT}$ against established VLM families, including BLIP-2 \citep{li2023blip}, InstructBLIP \citep{dai2023instructblip}, Qwen-VL \citep{bai2023qwen}, InternVL \citep{chen2024internvl}, Llava-1.5 \citep{vallaeys2024improved}, SmolVLM \citep{marafioti2025smolvlm}, PaLI \citep{chen2022pali}, PaliGemma \citep{beyer2024paligemma}, and Video-LLaVA \citep{lin2024video}. VL-JEPA$_\texttt{SFT}$ outperforms many of these baselines despite requiring significantly less computational resources--classical VLMs rely on extensively pretrained CLIP backbones combined with multi-stage visual instruction tuning. In comparison, VL-JEPA$_\texttt{SFT}$ employs a \textit{unified architecture} and a \textit{single embedding space} to seamlessly handle VQA, classification, and retrieval (Tab.~\ref{tab:cls_ret}).




\subsection{WorldPrediction-WM} 
\label{sec:exp.worldprediction}


\textbf{Evaluation Setup.} We evaluate VL-JEPA on the ``world modeling’’ task in the \textsc{WorldPrediction}~\citep{chen2025worldprediction} benchmark, where the model is provided with two images representing the initial and final world states and must identify, among four candidate video clips, the action that explains the observed transition. To adapt VL-JEPA, we duplicate and concatenate the initial and final state images to extract a \textit{state embedding}, and encode each action candidate into \textit{action embeddings}. The model then selects the candidate whose embedding is closest to the state embedding.

\textbf{Results.} Table~\ref{tab:worldprediction} shows accuracy comparisons. {VL-JEPA$_\texttt{BASE}$ attains \textbf{63.9\%} and VL-JEPA$_\texttt{SFT}$ attains \textbf{65.7\%}} top-1 accuracy on \textsc{WorldPrediction}-WM, establishing a new state of the art. Our VL-JEPA model not only substantially surpasses existing VLMs of comparable or larger scale but also exceeds the performance of frontier LLMs such as GPT-4o, Claude-3.5-sonnet, and Gemini-2.0. 

\subsection{Action Anticipation}

\begin{table}
    \centering
    \caption{Finetuning results on EPIC-KITCHENS-100 action anticipation (Recall@5) for different anticipation time $t$. VL-JEPA outperforms V-JEPA2 with a ViT-L-256px vision encoder across all anticipation time.}
    \resizebox{\linewidth}{!}{%~
    \begin{tabular}{cccccc}
    \toprule
         % &  & \multicolumn{4}{c}{Anticipation Time} \\\midrule
        \textbf{Model} & \textbf{Vision Encoder} & $t=1s$ & $t=2s$ & $t=4s$ & $t=10s$ \\\midrule
        
        {\color{mygray}V-JEPA2} & {\color{mygray}\texttt{ViT-g-384px}}  &{\color{mygray}\textbf{\underline{39.7}}}  & {\color{mygray}\textbf{\underline{28.6}}}  & {\color{mygray}19.3}  & {\color{mygray}8.2} \\
        % 
        V-JEPA2 & \texttt{ViT-L-256px} &  32.7 & 23.4 & 15.9 & 7.1\\%\midrule
        \rowcolor{headerblue}
        % 
        VL-JEPA & \texttt{ViT-L-256px} & 34.2 & 26.0 & \textbf{\underline{19.4}} & \textbf{\underline{11.7}}\\\bottomrule
        % VL-JEPA (w/o post-train)& 24.00 & 17.46 & 11.59 & 7.20\\
    \end{tabular}}
    \vspace{-10pt}
    \label{tab:action_anticipation}
\end{table}

\begin{table}[]
\centering
    \caption{Finetuning result on COIN next step forecasting. VL-JEPA established a new SoTA at 56.2.}
 \resizebox{\linewidth}{!}{%
    \begin{tabular}{ccccc}
    \toprule
    \textbf{Model} & \begin{tabular}[c]{@{}c@{}}VideoLLM\\ -online (8B)\end{tabular} & \begin{tabular}[c]{@{}c@{}}VideoLLM\\ -MoD (8B)\end{tabular} & \begin{tabular}[c]{@{}c@{}}ProVideLLM\\ -8B/11+\end{tabular} & \begin{tabular}[c]{@{}c@{}}VL-JEPA\\ (1.6B)\end{tabular}\\ 
        \midrule
    \begin{tabular}[c]{@{}c@{}}\textbf{Accuracy} \end{tabular} & 49.1 & 49.7 & 53.6 & \cellcolor{headerblue}\textbf{56.2} \\ \bottomrule
    \end{tabular}
}
\vspace{-10pt}
\label{tab:coin_forecasting}
\end{table}




\begin{figure*}
    \centering
    % \vspace{-5pt}
    \includegraphics[width=\linewidth]{figures/fair_comparison.png}
    % \vspace{-5pt}
    \caption{\textbf{Comparison of embedding prediction (VL-JEPA) and token prediction (VLM).} We conduct a fair comparison of under strictly aligned training settings (encoder, data, batchsize, etc.). \textbf{Left:} Zero-shot video captioning CIDEr score averaged over 3 datasets and zero-shot classification accuracy (top-5) averaged over 3 benchmarks. \textbf{Right:} Comparing the trainable parameters and average inference time cost.}
    \label{fig:vljepa_vs_vlm}
\end{figure*}



\textbf{Evauation Setup.} We assess the action anticipation (\textit{i.e.,} forecasting future actions) capabilities of fine-tuned VL-JEPA on two benchmarks: action anticipation with EPIC-KITCHENS-100 \citep{damen2022rescaling}, and COIN \citep{tang2019coin} next step forecasting. 
% EK100 is a widely adopted dataset for egocentric human-action anticipation and comprises approximately 100 hours of first-person cooking activity videos, providing a challenging testbed for temporal action prediction. 
In this task, given a window of context video, the model is required to predict the next action that will occur after a specified anticipation time—the temporal gap between the end of the context segment and the onset of the subsequent action. 
% While the default anticipation interval is set to 1 second, we additionally evaluate model performance across a range of anticipation times: 1s, 2s, 4s, and 10s. Performance is measured using mean class recall at rank 5 (Recall@5), reflecting the multi-class nature of the problem. For all experiments, 
VL-JEPA$_\texttt{SFT}$ is fine-tuned and we compare our results against the state-of-the-art baselines. 
% V-JEPA2 models \citep{assran2025v}, including V-JEPA2 ViT-L and V-JEPA 2 ViT-g384, as well as the VJEPA with probe variant, which currently achieves the highest reported performance on this task. Additionally, another benchmark we evaluate on is COIN next step forecasting. Given the previous segements of steps on procedural plan, the model needs to predict next step. Here, the anticipation time varies depending on samples and we follow the setup from \citep{lin2022learning} and compare performance with SOTA results of ProVideLLM \citep{chatterjee2025memory} by top-1 accuracy following previous work.  

\textbf{Results} As compated in Table VL-JEPA demonstrates strong performance in both benchmarks (Table~\ref{tab:action_anticipation} and Table~\ref{tab:coin_forecasting}). For the Epic-Kitchens-100, particularly at the standard 1-second anticipation interval, where it achieves a Recall@5 of 34.18 -- surpassing the previous V-JEPA 2 that uses the same ViT-L-256px encoder by 1.48 points. While the larger V-JEPA 2 ViT-g-384px model attains the highest score at 1 second, VL-JEPA remains competitive and shows notable advantages at longer anticipation intervals. 
% Specifically, VL-JEPA outperforms V-JEPA 2 ViT-g$_{384}$ at both 4 seconds (19.43 vs. 19.30) and 10 seconds (11.71 vs. 8.20), indicating its robustness in predicting future actions over extended time horizons.
On the COIN next step forecasting task, as shown in Tab. \ref{tab:coin_forecasting}, VL-JEPA achieves 56.2\%, outperform strong VLM baselines VideoLLM-online \citep{chen2024videollm}, VideoLLM-MoD \citep{wu2024videollm}, ProVideLLM \citep{chatterjee2025memory}. Overall, these findings demonstrated that VL-JEPA can be finetuned to handle semantic uncertain prediction tasks.


\subsection{Embedding Prediction vs. Token Prediction: A Controlled Comparison}
\label{sec:exp.embbeddings_vs_tokens}

\textbf{Evaluation Setup.} In this section, we compare VL-JEPA to a token-generative VLM baseline under a strictly aligned training conditions. Both models use the same Perception Encoder \citep{bolya2025perception} (frozen ViT-L-14 with 336$^2$ resolution, no tiling, 16 frames per video) for vision inputs. We use the same training iterations with the same effective batch size of 128, same learning rate scheduler on the same pretraining data mixture described above (\S\ref{sec:implementation-details}). The only difference is the prediction task: VL-JEPA predicts target embeddings \citep{duquenne2023sonar} using a 0.5B predictor, whereas the VLM baseline performs next-token prediction with cross-entropy using a 1B LLM. For VLM, we use the standard training recipe and codebase of PerceptionLM \citep{cho2025perceptionlm}, aligning frozen vision encoder and text-only LLM \texttt{Llama-3.2-1B}. For VL-JEPA, we initialize the predictor from the 8-16 layers of \texttt{Llama-3.2-1B}.

We evaluate both models at regular checkpoints throughout training spanning from 500K to 15M samples seen. At each checkpoint, we measure the performance on video captioning and video classification. For video captioning, we report CIDEr scores averaged across YouCook2 \citep{zhou2018towards}, MSR-VTT \citep{xu2016msr} and PVD-Bench \citep{bolya2025perception}. VL-JEPA decodes the predicted embeddings while VLM generates the tokens directly. For video classification, we report top-5 accuracy averaged across CrossTask-Step, CrossTask-Task \citep{zhukov2019cross} and EgoExo4D \citep{grauman2024ego}. For VL-JEPA we choose the candidate with lowest cosine distance to the predicted embedding, while for VLM we pick the class with lowest perplexity.


\textbf{Results.} As shown in Fig.~\ref{fig:vljepa_vs_vlm}, both models yield comparable performance after 500K samples seen in both tasks, with respectively 1.23 and 1.35 CIDEr in video captioning and 14.9\% and 14.0\% top-5 accuracy for VL-JEPA and VLM. After a few iterations, we show that VL-JEPA's performance increase is much sharper compared to VLM, reaching 14.7 CIDEr and 35.3\% top-5 accuracy after 5M samples seen. This gap remains constant as training scales at 15M samples with 14.8 CIDEr and 41.0\% top-5 accuracy for VL-JEPA, while the VLM baseline yield respectively 7.1 CIDEr and 27.2\% top-5 accuracy. This controlled comparison highlights the benefit of predicting embeddings rather than tokens, showing both higher sample efficiency and stronger absolute performance.

We compare the inference cost of the above VL-JEPA and the VLM by pre-loading 64 video frames into memory and repeatedly decoding text 100 times with the same prompt, measuring the average time per sample. As shown in Fig.~\ref{fig:vljepa_vs_vlm} (right most), both models exhibit comparable latency when generating text. What differentiates our model from classical VLM is the decoupling between the prompt processing (``Query Embedding'') and the video encoder (``Encoder + Predictor'') from the text generation module (``Decoder''). This allows us to only use the first part of the model to perform retrieval and decode text only when needed (see Section~\ref{sec:exp.selective decoding} below), making our model more scalable for online video inference.




\subsection{Effectiveness of Selective Decoding}
\label{sec:exp.selective decoding}

\begin{figure*}
    \centering
    \includegraphics[width=1\linewidth]{figures/selective_decoding.png}
    \caption{
        \textbf{Evaluation of selective decoding.}
        \textbf{Left:} 
            We compare uniform sampling of decoding points at fixed intervals (red) and embedding-guided selective decoding (blue). Performance is measured by the average CIDEr score between each annotation $y$ and its closest decoded output $\hat{y}$. 
        \textbf{Right:} 
            Results on EgoExo4D show that selective decoding achieves a Pareto improvement over uniform sampling: for the same performance level, it requires fewer decoding operations.
    }
    \label{fig:selective_decoding}
\end{figure*}


\begin{table*}[t!]
\centering
\caption{{\textbf{Comparison of text-encoders performance}. We report triplet-based accuracy (\%) on SugarCrepe++ and VISLA datasets.}}
\label{tab:y-decoder}
\resizebox{\linewidth}{!}{%
\begin{tabular}{llr r|c| ccccc|c| cc}
\multicolumn{2}{c}{\multirow{3}{*}{\textbf{Model}}} & \multirow{3}{*}{\makecell{\textbf{\# Params.}\\(total)}} & \multirow{3}{*}{\makecell{\textbf{\# Params.}\\(text encoder)}} & \multicolumn{6}{c|}{\textbf{SugarCrepe++} \citep{10.5555/3737916.3738487}} & \multicolumn{3}{c}{\textbf{VISLA} \citep{dumpala2024vislabenchmarkevaluatingembedding}} \\
 & & & & \multirow{2}{*}{\textbf{Average}} & \makecell{Replace\\Attribute} & \makecell{Replace\\Object} & \makecell{Replace\\Relation} & \makecell{Swap\\Attribute} & \makecell{Swap\\Object} & \multirow{2}{*}{\textbf{Average}} & \multirow{2}{*}{Generic} & \multirow{2}{*}{Spatial} \\
 % & & & & & & & & & & & & \\
\toprule
CLIP & ViT-L & 389M & 85M & 44.5 & 56.7 & 83.0 & 42.5 & 27.0 & 13.5 & 34.5 & 37.6 & 31.3 \\
SigLIP2 & ViT-g & 1.9B & 708M & 56.5 & 66.9 & 74.4 & \underline{52.1} & \underline{58.4} & \underline{30.6} & \underline{40.4} & \underline{48.7} & 32.1 \\
PE-Core & ViT-G & 2.3B & 537M & \underline{58.6} & \textbf{73.6} & 90.6 & 48.9 & 53.2 & 26.5 & 38.3 & 45.2 & 31.4 \\
\midrule
\rowcolor[HTML]{EBF3FF}
\textbf{VL-JEPA$_\texttt{BASE}$} & ViT-L & 1.6B & 300M & \textbf{63.9} & \underline{72.2} & \underline{90.1} & \textbf{52.2} & \textbf{62.9} & \textbf{42.0} & \textbf{42.9} & \textbf{49.8} & \textbf{35.9} \\
\rowcolor[HTML]{EBF3FF}
\textbf{VL-JEPA$_\texttt{SFT}$} & ViT-L & 1.6B & 300M & 58.4 & 68.5 & \textbf{90.9} & 47.4 & 55.4 & 29.8 & 39.5 & 44.8 & \underline{34.2} \\
\bottomrule
\end{tabular}
}
\end{table*}


\looseness=-1
\textbf{Evaluation Setup.} We evaluate the effectiveness of VL-JEPA’s embedding-guided selective decoding on long-form video streams.  To this end, we design a benchmark task where the goal is to recover a temporal sequence of annotations while minimizing the number of text decoding operations, which dominate inference cost. As shown in Fig.~\ref{fig:selective_decoding} (left), decoding is performed only at selected points along the VL-JEPA embedding stream, yielding a sequence of $N$ decoded outputs $[(\hat{t}_1, \hat{y}_1), (\hat{t}_2, \hat{y}_2), \ldots, (\hat{t}_N, \hat{y}_N)]$. Each ground-truth annotation $[(t_1, y_1), (t_2, y_2), \ldots, (t_T, y_T)]$ is then aligned to its nearest decoded output in time (illustrated as $\circ \cdot \cdot \cdot \circ$ in Fig.~\ref{fig:selective_decoding}), and CIDEr is computed between matched pairs. We use the EgoExo4D \citep{grauman2024ego} validation set in procedural activity domains, which consists of 218 videos with an average duration of 6 minutes and about $T=143$ atomic action annotations per video.


As a baseline, we consider \textit{uniform sampling}, where decoding points are placed at fixed intervals regardless of the underlying video content. Standard streaming VLMs are limited to this strategy, whereas VL-JEPA supports a more effective alternative: \textit{adaptive selection} of decoding points guided by its predicted embeddings. We apply agglomerative clustering with temporal connectivity constraints \citep{murtagh2012algorithms} to partition the embedding sequence into $N$ segments of high intra-segment monosemanticity \citep{chen2024subobject}, measured by variance (\textit{i.e.,} Ward distance). The intuition is that within a semantically coherent segment, decoded outputs are highly similar, so decoding once per segment captures the essential information while greatly reducing overall decoding cost. The midpoint of each segment is then chosen as the decoding point, and decoding is performed either from the exact embedding or from the average-pooled embedding within the segment. 


\textbf{Results.} As shown in Fig.~\ref{fig:selective_decoding} (right), we sweep the average decoding frequency from 2.0 Hz down to 0.01 Hz (\textit{i.e.,} average intervals between consecutive decoding operations from 0.5s to 100s) by adjusting either the stride of uniform sampling or the number of clusters in adaptive selection. Across the entire range, adaptive selection consistently Pareto-dominates uniform sampling. In particular, selective decoding at 0.35 Hz (\textit{i.e.,} $\sim$2.85s interval) matches the performance of uniform decoding at 1 Hz, reducing decoding cost by $\sim$2.85$\times$. We further observe that average pooling provides consistent gains for both strategies, since it provides denoising and stabilization on embeddings prior feeding into the decoder.




\subsection{Evaluation of Y-Encoder}
\label{sec:exp.y-encoder}


\textbf{Evaluation Setup.} We evaluate whether the JEPA architecture improves the \texttt{Y-Encoder} by following the uni-modal text-only (TOT) evaluation setup.  
We use the hard-negative benchmarks SugarCrepe++ \citep{10.5555/3737916.3738487} and VISLA \citep{dumpala2024vislabenchmarkevaluatingembedding}.  
These datasets test sensitivity to semantic and lexical changes in image descriptions.  
Each dataset contains triplets: two semantically similar descriptions of the same image ($p1$ and $p2$), and one negative description ($n$) created by altering attributes, relations, or objects.  
We compare \texttt{Y-Encoders} from different models by computing the cosine similarity for all description pairs.  
We check that the similarity between positives $sim(p1,p2)$ is higher than both the similarity between each positive and the negative $sim(p1,n)$ and $sim(p2,n)$.  
We report accuracy (\%) across all samples.


\textbf{Results.} Table~\ref{tab:y-decoder} shows the performance of different models on text hard-negative benchmarks.
VL-JEPA$_\texttt{BASE}$ achieves a micro average accuracy of $63.9\%$ on SugarCrepe++ and $42.9\%$ on VISLA.
This is higher than the best other models: PE-Core scores $58.6\%$ on SugarCrepe++ and SigLIP2 scores $40.4\%$ on VISLA.
The finetuned VL-JEPA$_\texttt{SFT}$ model also achieves competitive results, with $58.4\%$ on SugarCrepe++ and $39.5\%$ on VISLA.
These results indicate that VL-JEPA$_\texttt{BASE}$ has a \texttt{Y-Encoder} that is more resilient to text hard-negatives.




\begin{table*}[t]
    \caption{\textbf{Ablation studies results}. The default setting adopted by VL-JEPA is marked in \colorbox{headerblue}{blue}. We calculate $\pm$delta within each group of ablations in comparison with the default setting.}
    \label{tab:ablations}
    \centering
    % \setlength{\tabcolsep}{3pt} % Tighten spacing slightly to fit 7 columns
    
    % --- LEFT COLUMN (Groups a, b, c) ---
    \begin{minipage}[t]{0.455\textwidth}
        \resizebox{\linewidth}{!}{%
        % Column setup: Label (Right), Val (Center), Delta (Left), ...
        \begin{tabular}{r r l r l r l}
            \toprule
             & \multicolumn{2}{c}{\begin{tabular}{@{}c@{}}\textbf{Classification}\\(Accuracy)\end{tabular}} & \multicolumn{2}{c}{\begin{tabular}{@{}c@{}}\textbf{Retrieval}\\(Recall@1)\end{tabular}} & \multicolumn{2}{c}{\begin{tabular}{@{}c@{}}{\textbf{VQA}}\\(Accuracy)\end{tabular}} \\
            \midrule
            
            % VL-JEPA
            {\color{mygray} VL-JEPA$_{\texttt{SFT}}$} & {\color{mygray} 75.4} & & {\color{mygray} 63.8} & & {\color{mygray} 74.2} & \\
            \bottomrule \toprule
            
            % Group (a)
            \multicolumn{7}{l}{\textit{\color{mygray} (a) Effectiveness of pretraining stage on caption data}} \\
            \rowcolor{headerblue}
            w/ Pretraining & 49.0 & & 47.5 & & 46.1 & \\
            w/o Pretraining & 27.3 & \negval{-21.7} & 30.2 & \negval{-17.3} & 42.5 & \negval{-3.6} \\
            \bottomrule \toprule
            
            % Group (b)
            \multicolumn{7}{l}{\textit{\color{mygray} (b) Learning rate multiplier for Y-Encoder}} \\
            \rowcolor{headerblue}
            \textit{multiplier} = 0.05 & 27.3 & & 30.2 & & 42.5 & \\
            \textit{multiplier} = 1.00 & 23.7 & \negval{-3.6} & 28.8 & \negval{-1.4} & 40.7 & \negval{-1.8} \\
            \textit{multiplier} = 0.10 & 26.9 & \negval{-0.4} & 30.2 & \negval{-0.0} & 42.9 & \posval{+0.4} \\
            \textit{multiplier} = 0.01 & 25.6 & \negval{-1.7} & 27.7 & \negval{-2.5} & 41.0 & \negval{-1.5} \\
            \textit{multiplier} = 0.00 & 20.0 & \negval{-7.3} & 25.9 & \negval{-4.3} & 41.4 & \negval{-1.1} \\
            \bottomrule \toprule
            
            % Group (c)
            \multicolumn{7}{l}{\textit{\color{mygray} (c) Loss function (with no projection head on top frozen text encoder)}} \\
            \rowcolor{headerblue}
            InfoNCE & 23.3 & & 30.3 & & 44.3 & \\
            Cosine & 16.5 & \negval{-6.8} & 20.2 & \negval{-10.1} & 46.6 & \posval{+2.3} \\
            L1 & 14.8 & \negval{-8.5} & 15.5 & \negval{-14.8} & 41.9 & \negval{-2.4} \\
            L2 & 13.5 & \negval{-9.8} & 11.7 & \negval{-18.6} & 43.7 & \negval{-0.6} \\
            \bottomrule
        \end{tabular}%
        }
    \end{minipage}
    \hfill
    % --- RIGHT COLUMN (Groups d, e) ---
    \begin{minipage}[t]{0.525\textwidth}
        \resizebox{\linewidth}{!}{%
        \begin{tabular}{r r l r l r l}
            \toprule
             & \multicolumn{2}{c}{\begin{tabular}{@{}c@{}}\textbf{Classification}\\(Accuracy)\end{tabular}} & \multicolumn{2}{c}{\begin{tabular}{@{}c@{}}\textbf{Retrieval}\\(Recall@1)\end{tabular}} & \multicolumn{2}{c}{\begin{tabular}{@{}c@{}}{\textbf{VQA}}\\(Accuracy)\end{tabular}} \\
            \midrule
            
            % Group (d)
            \multicolumn{7}{l}{\textit{\color{mygray} (d) Predictor architecture and initialization}} \\
            \rowcolor{headerblue}
            Layer 8-16 & 27.3 & & 30.2 & & 42.5 & \\
            Layer 0-2 & 24.3 & \negval{-3.0} & 27.8 & \negval{-2.4} & 40.1 & \negval{-2.4} \\
            Layer 0-4 & 25.1 & \negval{-2.2} & 28.9 & \negval{-1.3} & 43.6 & \posval{+1.1} \\
            Layer 0-8 & 27.2 & \negval{-0.1} & 29.3 & \negval{-0.9} & 43.4 & \posval{+0.9} \\
            Layer 0-16 & 27.4 & \posval{+0.1} & 31.0 & \posval{+0.8} & 45.5 & \posval{+3.0} \\
            % Layer 4-12 & 27.0 & \negval{-0.3} & 29.1 & \negval{-1.1} & 40.5 & \negval{-2.0} \\
            w/o Bi-direction Attention & 26.7 & \negval{-0.6} & 31.2 & \posval{+1.0} & 40.6 & \negval{-1.9} \\
            w/o Llama-3 Initialization & 28.1 & \posval{+0.8} & 30.4 & \posval{+0.2} & 40.6 & \negval{-1.9} \\
            \bottomrule \toprule
            
            % Group (e)
            \multicolumn{7}{l}{\textit{\color{mygray} (e) Y-Encoder (trainable linear projection on top of frozen text encoder)}} \\
            \rowcolor{headerblue}
            EmbeddingGemma-300M & 19.5 & & 24.1 & & 42.5 & \\
            Qwen3-Embedding-0.6B & 24.5 & \posval{+5.0} & 24.5 & \posval{+0.4} & 41.5 & \negval{-1.0} \\
            Qwen3-Embedding-4B & 27.7 & \posval{+8.2} & 26.6 & \posval{+2.5} & 38.1 & \negval{-4.4} \\
            Qwen3-Embedding-8B & 29.6 & \posval{+10.1} & 29.5 & \posval{+5.4} & 41.9 & \negval{-0.6} \\
            PE$_{\text{core}}$-B (356M) & 29.4 & \posval{+9.9} & 34.5 & \posval{+10.4} & 35.9 & \negval{-6.6} \\
            PE$_{\text{core}}$-L (356M) & 29.0 & \posval{+9.5} & 34.2 & \posval{+10.1} & 42.9 & \posval{+0.4} \\
            PE$_{\text{core}}$-G (539M) & 33.9 & \posval{+14.4} & 32.0 & \posval{+7.9} & 41.8 & \negval{-0.7} \\
            \bottomrule
        \end{tabular}%
        }
    \end{minipage}
\end{table*}




\subsection{Ablation Study}
\label{sec:exp.ablations}

\textbf{Evaluation Setup.} We study different design choices for VL-JEPA. Here we train all ablation models on the SFT stage data for 10K steps with a batch size of 512 (5M samples seen) and constant learning rate. We report average classification top-1 accuracy of 8 datasets (Tab.~\ref{tab:cls_ret}), average text-to-video retrieval recall@1 of 8 datasets (Tab.~\ref{tab:cls_ret}), and {average VQA accuracy of 4 datasets (CLEVR, GQA, TallyQA simple and complex). We report the results in Tab.~\ref{tab:ablations}.}

\textbf{Results.} 
{\textbf{(a) Pretraining.} 
    Dropping the first query-free pretraining stage on image and video captions significantly hurt performance, especially on classification (-21.7) and retrieval (-17.3). 
\textbf{(b) LR Multiplier.} 
    The sweet point of learning rate multiplier to the Y-Encoder is around 0.05 to 0.10. Either faster or slower learning degrades the performance.
\textbf{(c) Loss Function.} 
    InfoNCE generally give superior performance compared to cosine, L1, and L2 losses, with the only exception being cosine loss outperform InfoNCE on VQA. However, only InfoNCE has the anti-collapse regularization and can be applied with unfrozen Y-Encoder.}
\textbf{(d) Predictor.} 
    In terms of predictor size, more layers yield better performance, especially on VQA performance. We also see that if using the original causal attention instead of updating to bi-direction attention hurt VQA performance (-1.9), since query tokens are appended after visual tokens, and visual tokens are no longer able to attend to query tokens. Finally, we also see that LLama-3 initialization is beneficial to VQA performance, although vision-language alignment (classification and retrieval) is a bit worse compared to randomly initialized Transformer layers.
{\textbf{(e) Y-Encoder.} 
    We tried different text encoder as the Y-Encoder, and confirmed that VL-JEPA works well with other embedding models than EmbeddingGemma-300M. Generally, larger encoder leads to better performance, with visually aligned text encoders (PE models) has significant advantage in classification and retrieval.}


